% results.tex -- Enterprise Node Orchestrator Paper (RU-V5)

\section{Results}

\subsection{State Machine Validation}

All 46 test cases pass, validating the state machine implementation across
six categories: happy path cycle (6 transitions from Idle through Submitting
back to Idle), pipelined receiving (transaction acceptance during Proving
and Submitting states), error recovery (all operational states transition to
Error; Error recovers to Idle), invalid transition rejection (descriptive
errors for undefined transitions), history tracking (full audit trail with
metadata), and event emitter correctness (per-transition and per-state events).

\subsection{Orchestration Overhead}

Table~\ref{tab:overhead} presents per-phase latencies across all five
benchmark scenarios.

\begin{table*}[t]
\centering
\caption{Orchestration Benchmark Results (30--50 iterations per scenario, depth 32)}
\label{tab:overhead}
\begin{tabular}{lrrrrrr}
\toprule
\textbf{Scenario} & \textbf{Batch} & \textbf{Total E2E} & \textbf{Overhead} & \textbf{Proving} & \textbf{Overhead/E2E} & \textbf{Memory} \\
                   & \textbf{Size}  & \textbf{(ms)}      & \textbf{(ms)}     & \textbf{(ms)}    & \textbf{(\%)}         & \textbf{(MB)} \\
\midrule
overhead\_only     & 8  & 2,785  & 593  & 12$^\dagger$   & 21.3$^\ddagger$ & 84.8 \\
snarkjs d32 b8     & 8  & 15,537 & 594  & 12,765 & 3.8  & 85.6 \\
rapidsnark d32 b8  & 8  & 5,281  & 594  & 2,507  & 11.2 & 86.2 \\
snarkjs d32 b16    & 16 & 31,371 & 1,187 & 28,008 & 3.8  & 85.9 \\
rapidsnark d32 b64 & 64 & 18,889 & 4,704 & 12,006 & 24.9 & 87.1 \\
\bottomrule
\end{tabular}
\begin{flushleft}
\small $^\dagger$Timer overhead only (0~ms simulated proving). $^\ddagger$Excluding
proving time; overhead as fraction of non-proving E2E.
\end{flushleft}
\end{table*}

\subsubsection{Phase Breakdown}

Table~\ref{tab:phases} details the per-phase timings for the two primary
scenarios: batch-8 overhead isolation and the target batch-64 configuration.

\begin{table}[htbp]
\centering
\caption{Per-Phase Latency Breakdown}
\label{tab:phases}
\begin{tabular}{lrrrr}
\toprule
\textbf{Phase} & \multicolumn{2}{c}{\textbf{Batch 8}} & \multicolumn{2}{c}{\textbf{Batch 64}} \\
               & \textbf{Mean} & \textbf{P95} & \textbf{Mean} & \textbf{P95} \\
               & \textbf{(ms)} & \textbf{(ms)} & \textbf{(ms)} & \textbf{(ms)} \\
\midrule
Batch formation     & 11.39  & 12.67  & 90.31   & 92.43 \\
Witness generation  & 581.86 & 584.20 & 4,613.21 & 4,622.13 \\
Proving             & 12.31  & 15.98  & 12,006.26 & 12,016.35 \\
DAC attestation     & 171.83 & 174.23 & 171.89  & 174.97 \\
L1 submission       & 2,007.26 & 2,015.60 & 2,007.35 & 2,015.34 \\
Checkpoint          & $<$0.01 & $<$0.01 & $<$0.01 & $<$0.01 \\
\midrule
\textbf{Total}      & \textbf{2,784.8} & -- & \textbf{18,889.2} & -- \\
\bottomrule
\end{tabular}
\end{table}

\subsubsection{Overhead Scaling}

Orchestration overhead scales linearly with batch size at a rate of
approximately 73.5~ms per transaction, dominated by witness generation
(72~ms/tx). Batch formation contributes 1.4~ms/tx. The linear scaling
relationship is:

\begin{equation}
T_\text{overhead}(n) = 11.4 + 72.7 \cdot n \quad \text{(ms, } n = \text{batch size)}
\end{equation}

\noindent verified at $n \in \{8, 16, 64\}$: predicted values of 593, 1,175,
and 4,664~ms match measured values of 593, 1,187, and 4,704~ms within 1\%.

\subsubsection{Statistical Quality}

For the batch-8 overhead scenario (50 iterations): mean 2,784.8~ms, standard
deviation 7.8~ms, 95\% CI half-width 2.2~ms (0.08\% of mean). All scenarios
achieve CI half-width $<$1\% of mean, well below the 10\% acceptance
threshold.

\subsection{Pipeline Speedup}

Table~\ref{tab:pipeline} compares sequential and pipelined execution across
six configurations with 10 batches each.

\begin{table}[htbp]
\centering
\caption{Sequential vs. Pipelined Architecture (10 batches)}
\label{tab:pipeline}
\begin{tabular}{lrrrr}
\toprule
\textbf{Configuration} & \textbf{Seq.} & \textbf{Pipe.} & \textbf{Speed-} & \textbf{Tput} \\
                        & \textbf{(s)}  & \textbf{(s)}   & \textbf{up}     & \textbf{(tx/s)} \\
\midrule
snarkjs b8             & 155.1 & 149.8 & 1.04$\times$ & 0.53 \\
rapidsnark b8          & 52.5  & 47.2  & 1.11$\times$ & 1.69 \\
snarkjs b16            & 313.5 & 302.8 & 1.04$\times$ & 0.53 \\
rapidsnark b16         & 83.5  & 72.8  & 1.15$\times$ & 2.20 \\
snarkjs b64            & 1,569 & 1,526 & 1.03$\times$ & 0.42 \\
\textbf{rapidsnark b64} & \textbf{189.0} & \textbf{146.4} & \textbf{1.29}$\times$ & \textbf{4.37} \\
\bottomrule
\end{tabular}
\end{table}

\subsubsection{Speedup Analysis}

The pipeline speedup is a function of the ratio $\rho = T_\text{prep} /
T_\text{prove+submit}$. When preparation time is negligible relative to
proving ($\rho \to 0$), the pipeline offers no benefit (speedup $\to 1.0$).
As $\rho$ increases, more preparation work overlaps with the proving phase:

\begin{itemize}
\item snarkjs b8: $\rho = 591 / 14{,}920 = 0.04$, speedup 1.04$\times$
\item rapidsnark b8: $\rho = 591 / 4{,}670 = 0.13$, speedup 1.11$\times$
\item rapidsnark b64: $\rho = 4{,}733 / 14{,}170 = 0.33$, speedup 1.29$\times$
\end{itemize}

The maximum theoretical speedup (unlimited batch size, equal preparation
and proving time) approaches 2.0$\times$ but is constrained by the proving
phase floor.

\subsection{End-to-End Latency vs. Target}

Table~\ref{tab:e2e_target} evaluates each configuration against the
90-second target.

\begin{table}[htbp]
\centering
\caption{End-to-End Latency vs. 90-Second Target}
\label{tab:e2e_target}
\begin{tabular}{lrrll}
\toprule
\textbf{Backend} & \textbf{Batch} & \textbf{E2E} & \textbf{Under} & \textbf{Margin} \\
                  &                & \textbf{(s)} & \textbf{90s?}  & \\
\midrule
snarkjs   & 8  & 15.5  & Yes & 5.8$\times$ \\
rapidsnark & 8  & 5.3   & Yes & 17.0$\times$ \\
snarkjs   & 16 & 31.4  & Yes & 2.9$\times$ \\
rapidsnark & 16 & 9.2   & Yes & 9.8$\times$ \\
\textbf{snarkjs}   & \textbf{64} & \textbf{156.9} & \textbf{No} & \textbf{1.7$\times$ over} \\
rapidsnark & 64 & 18.9  & Yes & 4.8$\times$ \\
\bottomrule
\end{tabular}
\end{table}

The snarkjs WASM backend at batch size 64 exceeds the target by
1.7$\times$ (156.9~s vs. 90~s). All other configurations meet the target
with margins ranging from 2.9$\times$ to 17.0$\times$.

\subsection{Memory Footprint}

Heap memory usage is stable across all scenarios at 84.8--87.1~MB. The
2.3~MB variance corresponds to the additional transaction objects held in
memory for larger batch sizes. Production deployment with a real SMT
containing 100K entries is estimated at $\sim$320~MB (234~MB SMT from
RU-V1 + 85~MB orchestrator), within the 500~MB target.

\subsection{Benchmark Reconciliation}

Table~\ref{tab:reconcile} validates our measurements against published data.

\begin{table}[htbp]
\centering
\caption{Benchmark Reconciliation with Published Data}
\label{tab:reconcile}
\small
\begin{tabular}{@{}llll@{}}
\toprule
\textbf{Metric} & \textbf{Ours} & \textbf{Published} & \textbf{Assess.} \\
\midrule
Orch. overhead   & 593~ms  & push0: 5.2~ms   & Exp.$^*$ \\
Batch form.      & 11.4~ms & push0: 5.2~ms   & 2.2$\times$ \\
Mem./component   & 85~MB   & push0: 8--14~MB & Exp.$^\dagger$ \\
\bottomrule
\end{tabular}
\begin{flushleft}
\small $^*$push0 measures message routing only, not witness generation.
$^\dagger$push0 dispatchers are stateless; our node holds full SMT state.
\end{flushleft}
\end{table}

Batch formation latency (11.4~ms) is 2.2$\times$ the push0 P50 (5.2~ms),
consistent with our node performing 8 SHA-256 hash computations for mock
SMT inserts while push0 dispatchers perform message routing only. No metric
diverges by more than 10$\times$ from published benchmarks after accounting
for architectural differences.
